\chapter{一元实函数}\label{ch:b2:realfun}

分析转向函数列（见\cref{ch:b2:funcseq}）之前，
值得比本科一年级所要求的更细致地了解一元函数的概貌：
单调函数的间断能有多么复杂，
凸函数必定有多么好的正则性，以及导数
所享有的特殊性质（达布）。这些结构性结果
简洁精辟，深受出题者青睐。

\section{单调函数}

\begin{theorem}[单调函数的正则性]\label{thm:b2:realfun:monotone}
设 $f \colon I \to \R$ 在区间上单调递增。
\begin{enumerate}
  \item 在每个内点 $a$ 处，单侧极限存在：
        \[
        f(a^-) = \sup_{x < a} f(x) \;\leq\; f(a) \;\leq\;
        f(a^+) = \inf_{x > a} f(x) ;
        \]
        每个间断点都是一个\emph{跳跃}。
  \item $f$ 的间断点全体构成的集合至多
        可数。\index{单调函数!间断点}
\end{enumerate}
\end{theorem}

\begin{proof}
(1) 集合 $\{f(x) : x < a\}$ 非空，且以 $f(a)$ 为上界：
其上确界 $s$ 满足 $f(x) \to s$（当 $x \to a^-$ 时；任给
$\varepsilon$，存在 $f(x_0) > s - \varepsilon$，由单调性知
$f(x) \in \intoc{s - \varepsilon}{s}$，对 $x \in
\intoo{x_0}{a}$ 成立）。右侧同理。

(2) 对每个间断点 $a$，取一个非空开区间
$J_a = \intoo{f(a^-)}{f(a^+)}$（确为一个跳跃）。对间断点 $a < b$，
$J_a$ 与 $J_b$ 不相交：$f(a^+) \leq f(c)
\leq f(b^-)$（$c$ 为介于其间的任意点）。每个 $J_a$ 都含有一个有理数；
不同的间断点对应不同的有理数，于是得到间断点集到
$\Q$ 的一个单射，而后者可数
（见\cref{prop:b2:structures:countablestable}）。
\end{proof}

\begin{example}
此界是最佳的：取 $(r_n)$ 为 $\Q \cap
\intoo{0}{1}$ 的一个枚举，并令 $f(x) = \sum_{n : r_n \leq x} 2^{-n}$
（\omterm{def:b2:series:summable}{可和族}的定义，
见\cref{def:b2:series:summable}）。于是 $f$ 在
$\intcc{0}{1}$ 上单调递增，且恰在
$\intoo{0}{1}$ 中的每个有理数处不\omterm{def:b2:metric:continuity}{连续}（跳跃量为 $2^{-n}$，位于 $r_n$ 处）：单调函数
\emph{可以}在稠密的\omterm{def:b2:structures:countable}{可数集}上不\omterm{def:b2:metric:continuity}{连续}。
\end{example}

\begin{example}[跳跃量不超过上升量]\label{ex:b2:realfun:jumpbudget}
设 $f$ 在 $\intcc{a}{b}$ 上单调递增，则跳跃量受总预算的约束：
若 $a < c_1 < \dots < c_m < b$ 为间断点，其跳跃量分别为
$s_i = f(c_i^+) - f(c_i^-) > 0$，则取相互穿插的分点
$a < c_1 < t_1 < c_2 < \dots$，并
在每段上运用单调性，
\[
\sum_{i=1}^{m} s_i \;\leq\; f(b) - f(a) :
\]
总上升量控制着总跳跃量。由此推得：对
每个 $k$，至多有 $k\,\bigl(f(b) - f(a)\bigr)$ 个间断点，
其跳跃量 $\geq \frac1k$——这是对
\cref{thm:b2:realfun:monotone}~(2) 的定量加强，因为间断点集是
关于 $k$ 的这些有限集的可数并。在上面的有理跳跃函数上，
预算恰好用完：
跳跃量 $2^{-n}$ 之和为 $1 = f(1^+) - f(0^-)$，
这在显然的推广意义下成立。单调函数可以稠密地跳跃，
但只能在严格的预算之内。
\end{example}

\section{凸函数}

\begin{lemma}[斜率不等式]\label{lem:b2:realfun:slopes}
设 $f$ 为 $I$ 上的凸函数，$x < y < z$ 为 $I$ 中的点。则
\[
\frac{f(y) - f(x)}{y - x}
\;\leq\; \frac{f(z) - f(x)}{z - x}
\;\leq\; \frac{f(z) - f(y)}{z - y} :
\]
弦的斜率关于两个端点都递增。\index{凸函数
!斜率不等式}
\end{lemma}

\begin{proof}
将 $y = \frac{z - y}{z - x}\,x + \frac{y - x}{z - x}\,z$ 写成一个
凸组合，因两系数均为正且和为
$1$。由凸性得
\[
f(y) \;\leq\; \frac{z-y}{z-x}\,f(x) + \frac{y-x}{z-x}\,f(z).
\]
对左侧不等式，两边同减 $f(x)$，利用
$\frac{z-y}{z-x} - 1 = -\frac{y-x}{z-x}$:
\[
f(y) - f(x) \leq \frac{y - x}{z - x}\bigl(f(z) - f(x)\bigr),
\]
再除以 $y - x > 0$。对右侧不等式，
改从 $f(z)$ 中减去：
\[
f(z) - f(y) \geq f(z) - \frac{z-y}{z-x}f(x) -
\frac{y-x}{z-x}f(z)
= \frac{z - y}{z - x}\bigl(f(z) - f(x)\bigr),
\]
再除以 $z - y > 0$。上面两步本质上是同一个重心恒等式，
只是对照的端点不同。
\end{proof}

\begin{theorem}[凸函数的正则性]\label{thm:b2:realfun:convexreg}
设 $f$ 为区间 $I$ 上的凸函数。
\begin{enumerate}
  \item 在每个内点处，$f$ 具有有限的单侧
        导数 $f'_g \leq f'_d$；二者都是该点的递增函数，
        特别地，$f$ 在
        $I$ 的内部\omterm{def:b2:metric:continuity}{连续}（但在端点处未必\omterm{def:b2:metric:continuity}{连续}）。
  \item $f$ 位于每条\emph{支撑线}之上：对内点 $a$，
        及任意 $m \in \intcc{f'_g(a)}{f'_d(a)}$，
        \[
        f(x) \geq f(a) + m(x - a) \qquad (x \in I).
        \]
  \item （加权琴生不等式）对 $x_i \in I$ 及非负权 $\lambda_i
        \geq 0$, $\sum\lambda_i = 1$:
        \[
        f\Bigl(\sum_i \lambda_i x_i\Bigr) \leq \sum_i \lambda_i
        f(x_i) .\index{詹森不等式}
        \]
\end{enumerate}
\end{theorem}

\begin{proof}
(1) 固定内点 $a$。由\cref{lem:b2:realfun:slopes}，斜率
$\tau(h) = \frac{f(a + h) - f(a)}{h}$ 是
$h$ 的递增函数（两侧均成立，且 $\tau(h_-) \leq \tau(h_+)$，其中 $h_- < 0 <
h_+$）。因此 $\tau$ 当 $h \to 0^-$ 时有有限极限（递增，
且以任一右斜率为上界）——此即 $f'_g(a)$——并且当
$h \to 0^+$ 时极限为 $f'_d(a)$，且 $f'_g(a) \leq f'_d(a)$。
有限单侧导数蕴含 $a$ 处\omterm{def:b2:metric:continuity}{连续}。关于单调性：
对 $a < b$ 的内点，$f'_d(a) \leq \frac{f(b) - f(a)}{b -
a} \leq f'_g(b)$，同样由斜率不等式得出。

(2) 当 $x > a$ 时：$\frac{f(x) - f(a)}{x - a} \geq f'_d(a) \geq m$；
当 $x < a$ 时：$\frac{f(a) - f(x)}{a - x} \leq f'_g(a) \leq m$。二者
整理后即得所需结论。

(3) 对点数作归纳，完全仿照
本科一年级分册中的做法（两点情形即为定义）——或一步完成：
在 $a = \sum\lambda_i x_i$ 处应用 (2)，并将支撑线不等式
在点 $x_i$ 处以权 $\lambda_i$ 加权平均：$\sum_i
\lambda_i f(x_i) \geq f(a) + m\sum_i\lambda_i(x_i - a) = f(a)$。
\end{proof}

\begin{omfigure}
\begin{tikzpicture}
\begin{axis}[omaxis, width=8.2cm, height=5.4cm,
    xmin=-2.4, xmax=2.5, ymin=-0.9, ymax=4.5,
    xtick={-1.5, 0.5, 2}, ytick=\empty]
  \addplot[omDef, very thick, domain=-2.1:2.15, samples=80]
    {x^2};
  \addplot[omThm, thick, domain=-1.5:2, samples=2]
    {0.5*x + 3};
  \addplot[omProp, thick, domain=-1.1:2.2, samples=2]
    {x - 0.25};
  \node[font=\small, omThm] at (axis cs:-1.05,3.15)
    {弦};
  \node[font=\small, omProp] at (axis cs:1.75,0.95)
    {支撑线};
  \node[font=\small, omDef] at (axis cs:-1.95,2.6) {$f$};
\end{axis}
\end{tikzpicture}

{\small 凸性一图尽览：在 $-1.5$ 与 $2$ 之间，
函数 $f(x) = x^2$ 的图像位于其弦（定义）之下，
并位于 $x = 0.5$ 处支撑线之上
（见\cref{thm:b2:realfun:convexreg}~(2)）——本章周末问题中的每一个不等式，
都是对这两种位置关系的
重新整理。}
\end{omfigure}

\begin{example}[端点处的不连续]
在 $\intcc{0}{1}$ 上，函数 $f(0) = 1$，$f(x) = 0$（当 $x > 0$ 时）
是凸函数，但在端点 $0$ 处不\omterm{def:b2:metric:continuity}{连续}。结论（1）
不可改进。
\end{example}

\begin{example}[角点与支撑线族]\label{ex:b2:realfun:corner}
对 $f(x) = \abs x$，取 $a = 0$：单侧导数为
$f'_g(0) = -1$ 与 $f'_d(0) = +1$，且
由\cref{thm:b2:realfun:convexreg}~(2)可知，
\emph{每一个}斜率 $m \in \intcc{-1}{1}$ 都给出一条支撑线：
\[
\abs x \geq m\,x \qquad (x \in \R,\ -1 \leq m \leq 1),
\]
每个不等式恰在一条半直线上或在 $0$ 处取等号。凸函数
在 $a$ 处可导，当且仅当支撑线族
退化为单一直线（$f'_g(a) = f'_d(a)$）；角点处则
有一区间的切线。此支撑线族正是凸优化中
\emph{次微分}的有限维雏形——
而凸函数之所以如此稳健，正因为即使导数失效，
支撑几何依然存在，这正是
琴生证明所用的全部。
\end{example}

\begin{example}[幂平均不等式]\label{ex:b2:realfun:powermeans}
设 $0 < p < q$，$x_i$ 为正数，权 $\lambda_i$
之和为 $1$，对凸函数 $t \mapsto t^{q/p}$
在点 $x_i^p$ 处应用琴生不等式，得
\[
\Bigl(\sum \lambda_i x_i^{p}\Bigr)^{1/p}
\leq \Bigl(\sum \lambda_i x_i^{q}\Bigr)^{1/q} :
\]
幂平均随指数增大——它包含算术–平方平均不等式不等式，并且
在极限 $p \to 0$（见\cref{exo:b2:realfun:6}）时，又得到
算术–几何平均不等式。
\end{example}

\begin{omfigure}
\begin{tikzpicture}
\begin{axis}[omaxis, width=8.6cm, height=5.6cm,
    xmin=-8.5, xmax=8.5, ymin=0.5, ymax=4.5,
    xtick={-4,-1,1,2,4}, ytick={1,2,4},
    xlabel={$p$}, ylabel={$M_p$}]
  \addplot[omDef, very thick, domain=-8:-0.4, samples=90]
    {((1/3)*(1^x + 2^x + 4^x))^(1/x)};
  \addplot[omDef, very thick, domain=0.4:8, samples=90]
    {((1/3)*(1^x + 2^x + 4^x))^(1/x)};
  \addplot[omThm, dashed, domain=-8.5:8.5, samples=2] {4};
  \addplot[omThm, dashed, domain=-8.5:8.5, samples=2] {1};
  \node[font=\small, omDef] at (axis cs:-5.4,1.55)
    {$M_p(1,2,4)$};
  \node[font=\small, omThm] at (axis cs:-6.4,3.75)
    {$\max = 4$};
  \node[font=\small, omThm] at (axis cs:6.3,1.25)
    {$\min = 1$};
  \node[font=\small, omProp] at (axis cs:0.2,1.83) {$2$};
\end{axis}
\end{tikzpicture}

{\small 幂平均 $M_p$ 作用于取值 $1, 2, 4$（等权），
作为指数 $p$ 的函数，单调递增：
自 $\min = 1$（$p \to -\infty$ 时）增至 $\max = 4$（$p \to
+\infty$ 时），依次经过调和平均（$p = -1$）、几何平均（
$p = 0$ 处的缺口，取值 $2$）、算术平均（$p = 1$）与平方平均（$p
= 2$）。整条经典均值不等式链
即为一条单调上升的曲线——在本章周末问题中
第三部分予以证明。}
\end{omfigure}

\begin{example}[最大熵]\label{ex:b2:realfun:entropy}
对概率向量 $(p_1, \dots, p_n)$（分量为正，
和为 $1$），熵 $H(p) = -\sum_i p_i\ln p_i$
满足
\[
H(p) \leq \ln n ,
\qquad\text{当且仅当 } p_i = \frac1n \text{ 对所有 }
i .
\]
证明：用詹森不等式（\cref{thm:b2:realfun:convexreg}~(3)），
应用于\emph{凹}函数 $\ln$，取权为 $p_i$、点为
$\frac{1}{p_i}$：
\[
H(p) = \sum_i p_i\ln\frac{1}{p_i}
\leq \ln\Bigl(\sum_i p_i\,\frac1{p_i}\Bigr) = \ln n ,
\]
取等号时，所有点 $\frac1{p_i}$ 必相等（严格
凹性），即 $p$ 为均匀分布。等价地，它即是
\cref{exo:b2:realfun:7}在 $q$ 为均匀分布时的情形。不确定性在无知均匀散布时达到最大——
这正是编码、统计力学，以及
在\cref{ch:b2:randomvar}中出现的熵背后的
变分原理。
\end{example}

\begin{method}[求隐藏在不等式背后的凸函数]\label{met:b2:realfun:spotting}
大多数经典不等式都是伪装了的詹森不等式；要将其还原，
可按下列步骤进行：(1) 归一化，使出现\emph{加权平均}（
权为正且和为 $1$——必要时除以总质量）；
(2) 观察平均号内外分别作用的是什么函数；
断言“$f(\text{平均}) \leq$
$f$ 的平均”时，$f$ 即为凸函数；(3) 用二阶导数验证凸性，
并通过严格性处理等号；(4) 若看不出平均，
先取对数——乘积与幂会变成平均，
而 $\ln$ 的凹性恰好给出算术–几何平均不等式、杨不等式及其同类
（本章周末问题对其中每一个都执行步骤 1--4）。若
取对数后仍看不出平均，
可尝试把不等式读作斜率的单调性
（\cref{lem:b2:realfun:slopes}）——
次可加性陈述（如\cref{exo:b2:realfun:9}）即属
此类。
\end{method}

\begin{remark}[常见误区]\label{rem:b2:realfun:pitfalls}
(i) 凸性不能被乘积保持：$x$ 与 $(x - 1)^2$
在 $\intcc{0}{2}$ 上都是凸的，但其乘积 $x(x-1)^2$ 的
二阶导数为 $6x - 4$，在
$\intco{0}{\frac23}$ 上为负，故它不是凸的；
凸性也不能在缺乏单调性时被复合保持
（\cref{exo:b2:realfun:10}）。(ii) 对凹函数，詹森不等式
的方向相反：一半的经典不等式是凹函数的
$\ln$ 版本；把凸形式用于 $\ln$，正是
\emph{反向}证明算术–几何平均不等式最快的办法。(iii) 仅由中点
凸性推不出凸性——还需要\omterm{def:b2:metric:continuity}{连续}性（或仅需有界性）
（\cref{exo:b2:realfun:8}）；
病态的反例超出了本书的公理范围。(iv)
\emph{开}区间上的凸函数是\omterm{def:b2:metric:continuity}{连续}的，
甚至是局部\omterm{def:b2:metric:continuity}{利普希茨}的（\cref{exo:b2:realfun:12}）；
在端点处则不然。(v) 导数满足达布性质，但未必
\omterm{def:b2:metric:continuity}{连续}（\cref{ex:b2:realfun:oscillation}）：
“$f'$ 没有跳跃”绝不意味着“$f'$ \omterm{def:b2:metric:continuity}{连续}”。
\end{remark}

\section{达布性质}

\begin{theorem}[达布]\label{thm:b2:realfun:darboux}
设 $f$ 在区间 $I$ 上可导，则 $f'$ 取遍
其任意两个取值之间的每一个值——尽管 $f'$ 未必
\omterm{def:b2:metric:continuity}{连续}。\index{达布定理}
\end{theorem}

\begin{proof}
设 $a < b$ 为 $I$ 中的点，$v$ 严格介于 $f'(a)$ 与 $f'(b)$ 之间，
不妨设 $f'(a) < v < f'(b)$。函数 $g(x) = f(x) - vx$
可导，且 $g'(a) < 0 < g'(b)$：其在
$\intcc{a}{b}$ 上的最小值（由紧集上的\omterm{def:b2:metric:continuity}{连续}性可取到）既不在 $a$ 处取得
（$a$ 之后 $g$ 减小至 $g(a)$ 以下），也不在 $b$ 处取得（
$b$ 之前 $g$ 低于 $g(b)$）：最小值点位于内部，且在该点处 $g'(c)
= 0$，即 $f'(c) = v$。（这曾是本科一年级的星号练习；其
在理论中的位置在此。）
\end{proof}

\begin{example}[哪些函数是导数？]\label{ex:b2:realfun:whichderivatives}
达布定理是一台“不存在性”的机器。取整函数
$\lfloor x\rfloor$ 不是 $\R$ 上任何函数的导数：
它取到值 $0$ 与 $1$，却跳过了 $\frac12$，
即在 $\intcc{0}{1}$ 上，而\cref{thm:b2:realfun:darboux}禁止导数这样做。
同样的结论适用于每一个带跳跃的函数，
例如符号函数、赫维赛德函数与一切阶梯函数，无论看起来多么无害；
它们的“原函数”（如符号函数对应的 $\abs x$ 等）
仅在跳跃点之外存在，并在该点处形成折角。
对比：极度不\omterm{def:b2:metric:continuity}{连续}的 $f'$（见\cref{ex:b2:realfun:oscillation}）
\emph{确实}是导数——其间断性属于振荡型，而达布性质
容许这种间断。
两类行为之间的分界，恰是
下面的无跳跃推论。
\end{example}

\begin{corollary}\label{cor:b2:realfun:nojumps}
导数没有跳跃型间断点：若 $f'(a^-)$ 与
$f'(a^+)$ 存在，则它们都等于 $f'(a)$。导数的间断点
总是振荡型的（例如 $x^2\sin\frac1x$ 的
导数在 $0$ 处，见本科一年级分册）。
\end{corollary}

\begin{proof}
若 $f'(a^+) = \lim_{x\to a^+} f'(x)$ 存在，且与
$f'(a)$ 不等，则 $f'$ 在右邻域内会跳过二者之间的值，
这与达布性质在区间
$\intcc{a}{a + h}$ 上的结论相矛盾。（另证：由中值定理可得
$f'(a) = \lim_{h\to0^+} \frac{f(a+h)-f(a)}{h} = f'(a^+)$，
因差商是某中间点处的 $f'$ 值。）
左侧同理。
\end{proof}

\begin{example}[典型的振荡导数]\label{ex:b2:realfun:oscillation}
设 $f(x) = x^2\sin\frac1x$（$x \neq 0$），且 $f(0) = 0$。
在 $0$ 处：$\bigl|\frac{f(h) - f(0)}{h}\bigr| = \abs{h\sin\frac1h}
\leq \abs h \to 0$，故 $f'(0) = 0$ 存在。在 $0$ 以外，
\[
f'(x) = 2x\sin\frac1x - \cos\frac1x ,
\]
其中第一项趋于 $0$，而 $\cos\frac1x$ 振荡
地取遍 $\intcc{-1}{1}$ 中的值，在每个区间 $\intoo{0}{\delta}$ 上
极限 $f'(0^+)$ 不存在。于是 $f'$ 处处有定义，但
在 $0$ 处不\omterm{def:b2:metric:continuity}{连续}——而且，正如
\cref{cor:b2:realfun:nojumps} 所预言的，这一间断是振荡，而非跳跃：
在每个 $\intoo{0}{\delta}$ 上，$f'$ 仍
扫过围绕 $0$ 的一个完整区间。导数可以极不规则，但
只能以与达布性质相容的方式出现。
\end{example}

\begin{remark}[本章的应用]
凸性是不等式这座“工厂”的引擎：本章的周末问题
由它推出杨、赫尔德与闵可夫斯基不等式，
以及由此得到的幂平均链，它们被\cref{ch:b2:nvs}的\omterm{def:b2:nvs:norm}{范数}理论
和\cref{ch:b2:integration}的积分估计所使用；
琴生不等式在概率论中以
\cref{ch:b2:randomvar}中的矩不等式的形式再次出现。单调正则性在
\cref{ch:b2:integration}中重现（单调函数可积），并且
在本科三年级分册中，以几乎处处可微性的形式
出现于单调函数——在那里，“可数多个跳跃”成为
勒贝格理论的第一步。
\end{remark}

\section{练习}

\begin{exercise}[$\star$]\label{exo:b2:realfun:1}
求下列函数的不\omterm{def:b2:metric:continuity}{连续}点集与跳跃大小：$\lfloor x
\rfloor$；$\;x - \lfloor x\rfloor$；$\;\lfloor x \rfloor + \sqrt{x
- \lfloor x\rfloor}$；
\cref{thm:b2:realfun:monotone}之后例中的函数，限制在二进有理数
$r_n$ 上。
\end{exercise}

\begin{exercise}[$\star$]\label{exo:b2:realfun:2}
证明：递增函数 $f \colon I \to \R$ 若具有介值性质
（即其在任一子区间上的像是区间），
则它是\omterm{def:b2:metric:continuity}{连续}的。
\end{exercise}

\begin{exercise}[$\star$]\label{exo:b2:realfun:3}
下列哪些函数在其定义域上是凸的？$x \mapsto
x\ln x$（$x > 0$）；$\;x \mapsto \ln(1 + \eu^x)$；$\;x \mapsto
\sqrt{1 + x^2}$；$\;x \mapsto x^3$。
\end{exercise}

\begin{exercise}[$\star\star$]\label{exo:b2:realfun:4}
设 $f$ 为 $\R$ 上的凸函数，且有上界。证明 $f$
为常值函数。\emph{（若 $f(a) \neq f(b)$，则斜率不等式
将非零弦斜率向外传播：越过取较大值的那一点后，
$f$ 至少线性增长——这与有界性矛盾。
斜率的正负两种情形都要讨论。）}由此推出：
$\R$ 上两端都有渐近线的凸函数是仿射函数。
\end{exercise}

\begin{exercise}[$\star\star$]\label{exo:b2:realfun:5}
设 $f$ 在 $I$ 上可微，且 $f'$ 单调。证明
$f'$ \omterm{def:b2:metric:continuity}{连续} \emph{（结合
\cref{thm:b2:realfun:monotone}与
\cref{cor:b2:realfun:nojumps}）}。
\end{exercise}

\begin{exercise}[$\star\star$]\label{exo:b2:realfun:6}
（几何平均作为极限）设 $x_i$ 为正数，权
$\lambda_i$ 之和为 $1$，证明
\[
\lim_{p \to 0^+} \Bigl(\sum_i \lambda_i x_i^p\Bigr)^{1/p}
= \prod_i x_i^{\lambda_i} ,
\]
利用 $x_i^p = \eu^{p\ln x_i} = 1 + p\ln x_i + O(p^2)$，
并由
\cref{ex:b2:realfun:powermeans}推出加权算术–几何平均不等式。
\end{exercise}

\begin{exercise}[$\star\star$]\label{exo:b2:realfun:7}
（熵不等式）利用 $t \mapsto t\ln t$ 的严格凸性，
证明：对正数 $p_i, q_i$，若 $\sum p_i = \sum q_i = 1$，则
\[
\sum_i p_i \ln\frac{p_i}{q_i} \geq 0 ,
\]
且当且仅当 $p = q$ 时等号成立。\emph{（将左边写成 $\sum q_i\,
\varphi\bigl(\frac{p_i}{q_i}\bigr)$ 的形式，其中 $\varphi(t) = t\ln t$，并
对权为 $q_i$ 的情形应用琴生不等式。）}
\end{exercise}

\begin{exercise}[$\star\star\star$]\label{exo:b2:realfun:8}
（中点凸性）$f \colon I \to \R$ 称为\emph{中点凸函数}，
当 $f\bigl(\frac{x+y}{2}\bigr) \leq \frac{f(x) + f(y)}{2}$
恒成立。证明：\emph{\omterm{def:b2:metric:continuity}{连续}}的中点凸函数是
凸的。\emph{（先对二进权
$\frac{k}{2^m}$ 用对 $m$ 的归纳法建立凸性不等式，
再利用稠密性与\omterm{def:b2:metric:continuity}{连续}性取极限。）}
\end{exercise}

\begin{exercise}[$\star\star\star$]\label{exo:b2:realfun:9}
设 $f$ 在 $\intco{0}{+\infty}$ 上为凸函数，且 $f(0) \leq 0$。证明
$x \mapsto \frac{f(x)}{x}$ 在
$\intoo{0}{+\infty}$ 上递增，并由此推出：对凸函数 $f$，若 $f(0) =
0$，则 $f(x + y) \geq f(x) + f(y)$，其中 $x, y \geq 0$
（超可加性）。
\end{exercise}

\begin{exercise}[$\star$]\label{exo:b2:realfun:10}
设 $f$ 为 $I$ 上的凸函数，$g$ 为
包含 $f(I)$ 的某个区间上的凸且\emph{递增}函数。证明 $g \circ f$ 为凸函数，
并用反例说明 $g$ 的单调性不能
省去。
\end{exercise}

\begin{exercise}[$\star\star$]\label{exo:b2:realfun:11}
（埃尔米特–阿达马）设 $f$ 为凸函数，且在
$\intcc{a}{b}$ 上\omterm{def:b2:metric:continuity}{连续}。证明
\[
f\Bigl(\frac{a+b}{2}\Bigr) \;\leq\; \frac{1}{b -
a}\int_a^b f(t)\,\dd t \;\leq\; \frac{f(a) + f(b)}{2} .
\]
\emph{（左：在中点处对支撑线积分；右：
用弦给 $f$ 作上界。）}
\end{exercise}

\begin{exercise}[$\star\star\star$]\label{exo:b2:realfun:12}
证明：定义在\emph{开}区间 $I$ 上的凸函数
是局部\omterm{def:b2:metric:continuity}{利普希茨}的：对每个线段 $\intcc{a}{b} \subseteq I$，
及裕量 $\delta > 0$，使得 $\intcc{a - \delta}{b + \delta}
\subseteq I$，则 $f$ 在 $\intcc{a}{b}$ 上的限制
为\omterm{def:b2:metric:continuity}{利普希茨}函数，其常数为
$\max\Bigl(\bigl|\frac{f(a) - f(a - \delta)}{\delta}\bigr|,
\bigl|\frac{f(b + \delta) - f(b)}{\delta}\bigr|\Bigr)$
\emph{（夹住这两者之间的每个弦斜率，
借助斜率不等式）}。
\end{exercise}

\section{周末问题：凸性工具箱}

一个定义——图像上方的弦——\omterm{def:b2:structures:generated}{生成}了
经典不等式的全部工具。本周末问题按逻辑顺序构建这套工具：
先是凸性判别准则与严格琴生不等式，
然后是杨、赫尔德与闵可夫斯基不等式（它们是 $p$-范数的奠基性结果），
以及从最小到最大的完整幂平均链，
还有两项额外成果——卡莱曼不等式，以及
将赫尔德不等式看作一种对偶。一切结论
均给出证明，不引入任何外部结论。

\begin{problem}[{周末问题——杨、赫尔德、闵可夫斯基，
以及幂平均链}]
\label{pb:b2:realfun:1}
全文，$p, q > 1$ 称为\emph{共轭指数}，满足 $\frac1p +
\frac1q = 1$；向量记作 $a = (a_1, \dots, a_n) \in \R^n$；
权 $\lambda_i > 0$ 满足 $\sum_i\lambda_i = 1$。

\medskip
\noindent\textbf{第一部分——判别准则与严格琴生不等式。}
\begin{enumerate}
  \item 设 $f$ 为区间 $I$ 上的可微函数。证明
        $f$ 为凸函数当且仅当 $f'$ 递增，
        \emph{（一个方向是在\cref{lem:b2:realfun:slopes}
        的斜率不等式中取极限；另一方向
        用中值定理）}。由此推出 $C^2$ 的
        判别准则 $f'' \geq 0$。
  \item 设 $f''> 0$ 在 $I$ 上成立。证明 $f$ 是
        \emph{严格}凸（对 $x \neq y$
        与 $\lambda \in \intoo01$ 时严格不等式成立），并证明严格凸函数
        满足琴生不等式（
        见\cref{thm:b2:realfun:convexreg}~(3)），等号
        \emph{仅当}所有 $x_i$ 相等时成立。
  \item 验证工具箱的基本材料：$-\ln$ 在
        $\intoo{0}{+\infty}$ 上严格凸；$t \mapsto t^r$ 在该区间上，当 $r > 1$
        时严格凸，当 $0 < r < 1$ 时严格凹；$\exp$ 在 $\R$
        上严格凸。
  \item （杨不等式）对 $a, b \geq 0$，证明
        \[
        ab \;\leq\; \frac{a^p}{p} + \frac{b^q}{q},
        \]
        当且仅当 $a^p = b^q$ 时取等号 \emph{（
        对 $\ln$ 的凹性应用于两点 $a^p, b^q$，
        以 $\frac1p, \frac1q$ 为权）}。
  \item 一行内重新推导加权算术–几何平均不等式，由
        $\ln$ 的凹性出发：
        \[
        \prod_i x_i^{\lambda_i} \leq \sum_i\lambda_ix_i \qquad
        (x_i > 0),
        \]
        并指出等号成立的情形，与
        \cref{exo:b2:realfun:6}中的极限路线作比较。
\end{enumerate}

\medskip
\noindent\textbf{第二部分——赫尔德与闵可夫斯基。}记
$\norm{a}_p = \bigl(\sum_i \abs{a_i}^p\bigr)^{1/p}$ 及
$\norm{a}_\infty = \max_i\abs{a_i}$。
\begin{enumerate}[resume]
  \item （赫尔德不等式）证明
        \[
        \sum_{i=1}^{n}\abs{a_ib_i} \;\leq\;
        \norm a_p\,\norm b_q ,
        \]
        等号成立当且仅当向量 $(\abs{a_i}^p)$ 与
        $(\abs{b_i}^q)$ 成比例\emph{（先归一化
        $\norm a_p = \norm b_q = 1$，并对每一项应用杨不等式）}。
  \item 指出特殊情形：$p = q = 2$
        （柯西–施瓦茨不等式），以及端点对 $(p, q) = (1,
        \infty)$：叙述并证明 $\sum\abs{a_ib_i} \leq
        \norm a_1\norm b_\infty$。
  \item （闵可夫斯基不等式）对 $p \geq 1$，证明
        \[
        \norm{a + b}_p \leq \norm a_p + \norm b_p
        \]
        \emph{（先写出 $\abs{a_i + b_i}^p \leq \abs{a_i +
        b_i}^{p-1}(\abs{a_i} + \abs{b_i})$，并
        对每个乘积应用赫尔德不等式）}。于是 $\norm\cdot_p$ 是
        $\R^n$ 上的\omterm{def:b2:nvs:norm}{范数}，对每个 $p \in \intco{1}{+\infty}$，
        从而完成\cref{ch:b2:nvs}中的图景。
  \item 积分形式：设 $f, g$ 在
        $\intcc{a}{b}$ 上\omterm{def:b2:metric:continuity}{连续}，叙述并证明赫尔德与闵可夫斯基不等式，
        对 $\norm f_p = \bigl(\int_a^b\abs
        f^p\bigr)^{1/p}$ \emph{（证明方法相同，
        积分的严格正性用于等号情形的
        讨论）}。
  \item 证明单调性 $\norm a_q \leq \norm a_p$，其中 $1
        \leq p \leq q$，极限 $\norm a_p \to
        \norm a_\infty$，当 $p \to \infty$ 时，以及
        带精确常数的反向比较：
        \[
        \norm a_p \leq n^{\frac1p - \frac1q}\,\norm a_q
        \]
        \emph{（对常向量应用赫尔德不等式）}。指出
        每个等号成立的向量。
  \item （插值）设 $1 \leq p < r < q$ 且 $\theta \in
        \intoo01$ 满足 $\frac1r = \frac\theta p +
        \frac{1-\theta}q$，证明
        \[
        \norm a_r \leq \norm a_p^{\theta}\,
        \norm a_q^{1-\theta}
        \]
        \emph{（应用指数为 $\frac{p}{\theta
        r}$ 与 $\frac{q}{(1-\theta)r}$ 的赫尔德不等式，作用于 $\abs{a_i}^{\theta
        r}\abs{a_i}^{(1-\theta)r}$）}。
\end{enumerate}

\medskip
\noindent\textbf{第三部分——幂平均链，完整版。}
设 $p \neq 0$，令 $M_p = \bigl(\sum_i\lambda_i
x_i^p\bigr)^{1/p}$（$x_i > 0$），并令 $M_0 = \prod_i
x_i^{\lambda_i}$。
\begin{enumerate}[resume]
  \item 证明 $p \mapsto M_p$ 在整个
        $\R^*$ 上递增：对 $p < q < 0$，由倒数恒等式
        $M_{-p}(x) = M_p(1/x)^{-1}$，再经由 $0$ 过渡，
        证明 $M_p \leq M_0 \leq M_q$ 对 $p < 0 < q$ 成立
        \emph{（应用 $\ln$ 的凹性于 $x_i^q$，并
        利用负指数下的反向不等式）}。
  \item 证明极限 $M_p \to \max_i x_i$（当 $p \to
        +\infty$ 时），以及 $M_p \to \min_i x_i$（当 $p \to -\infty$ 时）。
  \item 写出等权情形下的链 $\min \leq \mathrm{HM} \leq
        \mathrm{GM} \leq \mathrm{AM} \leq \mathrm{QM} \leq
        \max$，并证明一个经典推论：
        对正数 $a_1, \dots, a_n$，
        \[
        \Bigl(\sum_i a_i\Bigr)\Bigl(\sum_i\frac1{a_i}\Bigr)
        \geq n^2 .
        \]
  \item 将平均与\omterm{def:b2:nvs:norm}{范数}联系起来：对等权 $\lambda_i =
        \frac1n$ 时，有 $M_p(x) = n^{-1/p}\norm x_p$。
        协调两种单调性：平均随 $p$ 增大而\emph{增大}，
        \omterm{def:b2:nvs:norm}{范数}则随之\emph{减小}（第 10 题）——请用一句话
        说明因子 $n^{-1/p}$ 的作用。
  \item 确定第 14 题整条链上的等号成立情形
        （正权重）：链上任一处取等号，都迫使所有 $x_i$
        相等——严格凸性由此发挥作用。
\end{enumerate}

\medskip
\noindent\textbf{第四部分——红利。}
\begin{enumerate}[resume]
  \item （带可调参数的杨不等式）对 $a, b \geq 0$ 与 $\varepsilon >
        0$，证明
        \[
        ab \leq \varepsilon\,\frac{a^p}{p} +
        \varepsilon^{-q/p}\,\frac{b^q}{q},
        \]
        以及主要特例 $ab \leq \varepsilon a^2 +
        \frac{b^2}{4\varepsilon}$，这是分析中通篇使用的
        吸收技巧。
  \item （迈向卡莱曼不等式）设 $c_k = \frac{(k+1)^k}{k^{k-1}}$。
        证明裂项相消恒等式 $\prod_{k=1}^{n}c_k =
        (n+1)^n$，并对
        $c_ka_k$ 这些数运用算术--几何平均不等式，推出
        \[
        (a_1a_2\cdots a_n)^{1/n} \leq
        \frac{1}{n(n+1)}\sum_{k=1}^{n} c_k a_k
        \qquad (a_k > 0).
        \]
  \item （卡莱曼不等式）对 $n$ 求和，并交换
        求和次序（正项\omterm{def:b2:series:summable}{可和族}，见
        \cref{thm:b2:series:fubini}），并利用 $\sum_{n \geq
        k}\frac{1}{n(n+1)} = \frac1k$ 及 $c_k/k = \bigl(1 +
        \frac1k\bigr)^k < \eu$，从而证得：对每个
        各项为正的收敛级数 $\sum a_k$，有
        \[
        \sum_{n=1}^{\infty}(a_1a_2\cdots a_n)^{1/n}
        \;\leq\; \eu\sum_{k=1}^{\infty}a_k .
        \]
  \item 设 $f$ 在 $\intcc{0}{1}$ 上\omterm{def:b2:metric:continuity}{连续}且为正，
        证明
        \[
        \Bigl(\int_0^1 f\Bigr)\Bigl(\int_0^1\frac1f\Bigr)
        \geq 1,
        \]
        且等号成立当且仅当 $f$ 为常值
        \emph{（对 $\sqrt f\cdot\frac1{\sqrt
        f}$ 应用柯西--施瓦茨不等式）}。
  \item （单位球的几何）利用
        闵可夫斯基不等式的等号成立情形，证明当 $1 < p < \infty$ 时，
        $\norm\cdot_p$ 的单位球面不含线段（该\omterm{def:b2:nvs:norm}{范数}
        严格凸的含义见
        \cref{pb:b2:nvs:1}），而当 $p = 1$ 与 $p =
        \infty$ 时则含线段，并指出其中的平坦部分。
\end{enumerate}

\medskip
\noindent\textbf{第五部分——对偶与综合。}
\begin{enumerate}[resume]
  \item （赫尔德不等式的对偶形式）证明对每个 $a \in
        \R^n$，
        \[
        \norm a_p = \max_{\norm b_q \leq 1}\ \sum_i a_ib_i ,
        \]
        并显式给出一个取最大值的 $b$。（$p$ \omterm{def:b2:nvs:norm}{范数}
        是 $q$ \omterm{def:b2:nvs:norm}{范数}的对偶——有限维情形下
        $L^p$ 对偶的雏形。）
  \item （矩）设 $X$ 为随机变量，取有限多个正值
        $x_i$，其概率分别为
        $\lambda_i$。将第 12 题改写为：$r \mapsto
        \E[X^r]^{1/r}$ 递增——即矩的
        李雅普诺夫不等式，将在
        \cref{ch:b2:randomvar}中重新使用。
  \item 借助有名不等式求解，每小题两行：（i）对
        正数 $a, b, c$：$a^3 + b^3 + c^3 \geq
        \frac{(a+b+c)^3}{9}$；（ii）对正数 $x_1, \dots,
        x_n$：$\bigl(\sum_i\sqrt{x_i}\bigr)^2 \leq
        n\sum_i x_i$。
  \item （综合）用五句话画出这套工具的谱系：由弦的定义到斜率引理；
        由斜率到支撑线，再到琴生不等式；
        由 $\ln$ 的凹性到杨不等式，再到赫尔德不等式，
        再到闵可夫斯基不等式，再到 $p$ \omterm{def:b2:nvs:norm}{范数}；由琴生不等式到幂平均链，
        再到矩；由算术--几何平均不等式到卡莱曼不等式。点明两座高峰
        （赫尔德--闵可夫斯基；卡莱曼），并指出这套工具的去向：
        本科三年级分册中的 $L^p$ 空间，
        其公理恰为第 6 题与第 8 题。
\end{enumerate}
\end{problem}
